\section{Prozessschritt 2: Risikobehandlung} \label{sec:methode2}
\subsection{Grundlegender Prozess und Strategien}
In der Praxis existiert keine einzelne standardisierte Risikominimierungsstrategie für Data-Science-Projekte. 
Stattdessen setzen Unternehmen auf eine Kombination verschiedener Maßnahmen. 
Dazu zählen insbesondere iterative Vorgehensweisen, bei denen Fehler frühzeitig erkannt und behoben werden, 
sowie die Anwendung international anerkannter Standards zur Risikoidentifikation und zum Risikomanagement. 
Mit Blick auf die in der Literatur identifizierten Risiken, beispielsweise mangelnde Team- und Data-Science-Kompetenzen, 
kommt zudem der HR-Abteilung eine wichtige Rolle beim Aufbau geeigneter Fachkompetenzen und interdisziplinärer Teams zu \cite{holtkemper_risk_2023}.

Die Vorgehensweise nach ISO 27001 fokussiert sich primär auf das Management von Informations- und IT-Sicherheitsrisiken, 
bietet jedoch aufgrund ihres strukturierten und präventiven Charakters eine sinnvolle Basis für das Risikomanagement 
in Data-Science-Projekten \cite{Kersten_27001_2019}. 

Entsprechend lässt sich der Prozess der Risikominimierung nach der ISO 27001
in folgende Schritte eingliedern:

\begin{itemize}
    \item Risikovermeidung
    \item Proaktiver Schutz
    \item Schadenbegrenzung
    \item Schaden- bzw. Risikoüberwälzung
    \item Akzeptanz des Restrisikos
\end{itemize}

Um Risiken entsprechend zu behandeln, können zunächst riskante Vorgehensweisen projektübergreifend bewusst vermieden oder entsprechend verändert werden,
 um als Organisation risikoarm agieren zu können. Systeme und Informationen können zudem proaktiv geschützt werden, 
 indem in Firewalls, Authentifizierungsprozesse, Zugriffskontrollen und Integritätsprüfungen investiert wird. Des Weiteren 
 können potenzielle Angriffe erschwert werden, indem auf die Verschlüsselung von Daten gesetzt oder betroffene Systeme an qualifizierte 
 Dienstleister ausgelagert werden. So wird das Risiko eines tatsächlich eintretenden Datendiebstahls deutlich reduziert. In Bezug auf Risiken
  mit potenziell hohem Schadenswert aber geringer Eintrittswahrscheinlichkeit können zusätzlich Versicherungen abgeschlossen werden 
  (z.B. Betriebsunterbrechungsversicherungen, Datenmissbrauchsversicherungen oder Versicherungen für Überspannungs-, Brand- oder Wasserschäden).
   Zu guter Letzt bleibt nach Durchführung aller bereits beschriebenen Safeguards jedoch ein Restrisiko bestehen, das nicht eliminiert werden kann. 
   Dieses muss so genau wie möglich bestimmt und laufend überwacht werden (vgl. \hyperref[sec:risiko_monitoring]{Kapitel~3.4}). Im Falle eines zunehmenden Restrisikos sind zusätzliche Maßnahmen einzuleiten.

   Klassische Safeguards zum Zwecke der IT-Sicherheit bilden im Rahmen von Data Science Projekten zwar ein solides Fundament greifen aber aufgrund der hohen Komplexität von Data Science Projekten zu kurz. Spezifische Risiken, wie "Bias", "Overfitting" oder "mangelnde Interpretierbarkeit der entwickelten Modelle"
   erfordern eine Evplution von einem statischen Risikomanagement hin zu einem automatisierten effizienten Risikomanagement durch MLOps. Hierfür eignen sich insbesondere Tools, wie "MLflow", "SHAP" oder "DQLabs". Diese dienen insbesondere der automatisierten Erkennung von dehlenden oder fehlerhaften Daten (DQLabs), 
   dem automatischen Tracking von Experimenten und der automatischen Identifikation von Bias und entsprechend verhinderten Freigabe eines Modells nach Übershcreiten bestimmter Schwellwerte (McKinsey9):

\subsection{Weitere Maßnahmen}

Parallel zur Definition und Durchführung relevanter Sicherheitsmaßnahmen stützen sich Unternehmen auf strukturierte Prozessmodelle wie CRISP-DM, um den Projektablauf zu standardisieren und Unsicherheiten zu reduzieren. Ein weiterer zentraler Ansatz ist die Sicherstellung hoher Datenqualität durch gezielte Datenvalidierung und -aufbereitung, was durch ein übergreifendes Data Governance Rahmenwerk gewährleistet werden kann. Hierfür ist es essenziell, dass Data Owners während der Dauer des gesamten Projekts eingebunden werden, um so die Entwicklung zu monitoren. Dadurch können Sie die Akzeptanz der Ergebnisse auf Seiten der Stakeholder erhöhen und die Effizienz des gesamten Prozesses durch hohe Datenqualität, effiziente Kommunikation und den einfachen Zugang zu Daten sicherstellen. Eine bestehende und stabile IT- sowie Dateninfrastruktur ist hierfür eine grundlegende Voraussetzung (Borus und Janssen, 2020). Interdisziplinäre Teams sowie die Orientierung an etablierten Risikomanagement-Frameworks sind in diesem Kontext zusätzliche ergebnisfördernde Maßnahmen (vgl. \hyperref[sec:begriffe]{Kapitel~2.1}).
Entsprechend des COSO ERM Framework spielt bezogen auf übergreifende Geschäftsrisiken vor allem die Unternehmenskultur und Governance eine übergeordnete Rolle. Dies lässt sich ebenfalls auf Risiken in Data Science Projekten übertragen. Durch den gezielten Ausbau von Führungsqualitäten, einer Vision, sowie einer Mission kann eine Kultur geschaffen werden, die ihre Mitarbeiter im Hinblick auf die nötige Agilität und Resilienz fördert und Akzeptanz für die entsprechenden Data Science Initiativen und ihre Ergebnisse schafft. Folglich wird ebenfalls das Erstellen eines Governace-Programms gefördert und die Mitarbeiter befähigt diesem Werk folge zu leisten COSO und KKotter \cite{COSO_2020}, \cite{Kotter_2011}. 

Audgrund der ethischen Sorgfaltspflicht im Rahmen von Data Science Projekten ist neben einem ausführliches Data Governance Rahmenwerk und einer fördernden Unternehmenskultur allerdings auch ein entsprechender Monitoring und Reporting-Prozess aufzusetzen, um Fehlverhalten und systematischen Bias zu verhindern (vgl. \hyperref[sec:risiko_monitoring]{Kapitel~3.3 Risikomonitoring}), \cite{Kersten_27001_2019.}

Alles in allem basiert erfolgreiche Risikominimierung in Data-Science-Projekten somit auf einem integrativen Ansatz, der verschiedene Methoden und die Anwendung mehrerer relevanter Standards vereint. Geeignete Normen sind hierfür ISO 31000, ISO/IEC 27001, ISO/IEC 25010 und ISO 21500. So können die verschiedenen Risiken, die in Data Science Projekten sowohl technischer, organisatorischer als auch gesetzlicher Natur sind, angemessen adressiert und minimiert werden. Dabei ist es der Organisation jederzeit überlassen auch eine eigens entwickelte Vorgehensweise anzuwenden, sie sollte lediglich regulatorische Anforderungen (DSGVO, EU-AI-Act etc.) erfüllen.